\documentclass[12pt,a4paper]{article}

% ─── Codificação e idioma ────────────────────────────────────────────────────
\usepackage[utf8]{inputenc}
\usepackage[T1]{fontenc}
\usepackage[brazil]{babel}

% ─── Fontes ─────────────────────────────────────────────────────────────────
\usepackage{lmodern}
\usepackage{microtype}

% ─── Matemática ─────────────────────────────────────────────────────────────
\usepackage{amsmath, amssymb, amsthm}
\usepackage{mathtools}
\usepackage{bm}

% ─── Layout e geometria ─────────────────────────────────────────────────────
\usepackage[
  top=2.5cm, bottom=2.5cm,
  left=2.5cm, right=2.5cm,
  headheight=15pt
]{geometry}

% ─── Cabeçalho/rodapé ───────────────────────────────────────────────────────
\usepackage{fancyhdr}
\pagestyle{fancy}
\fancyhf{}
\fancyhead[L]{\small\color{navyblue}\textbf{Intendência de Materiais --- Borba/AM}}
\fancyhead[R]{\small\color{navyblue}Versão 2.0 · 2025}
\fancyfoot[C]{\small\thepage}
\renewcommand{\headrulewidth}{0.6pt}

% ─── Cores ──────────────────────────────────────────────────────────────────
\usepackage[dvipsnames, table]{xcolor}
\definecolor{navyblue}{HTML}{1A3A5C}
\definecolor{medblue}{HTML}{2E75B6}
\definecolor{lightblue}{HTML}{D6E4F0}
\definecolor{rowalt}{HTML}{EBF3FB}
\definecolor{formulabg}{HTML}{F0F4FA}
\definecolor{alertbg}{HTML}{FFF8E1}
\definecolor{alertborder}{HTML}{F59E0B}

% ─── Tabelas ────────────────────────────────────────────────────────────────
\usepackage{booktabs}
\usepackage{array}
\usepackage{longtable}
\usepackage{multirow}
\usepackage{tabularx}
\usepackage{colortbl}
\renewcommand{\arraystretch}{1.4}

% ─── Caixas e destaques ─────────────────────────────────────────────────────
\usepackage{tcolorbox}
\tcbuselibrary{skins, breakable, theorems}

% caixa de fórmula
\newtcolorbox{formulabox}{
  enhanced, breakable,
  colback=formulabg, colframe=medblue,
  boxrule=0.8pt, arc=4pt,
  left=12pt, right=12pt, top=2pt, bottom=2pt
}

% caixa de alerta/correção
\newtcolorbox{alertbox}[1][Correção Técnica]{
  enhanced,
  colback=alertbg, colframe=alertborder,
  boxrule=1pt, arc=4pt,
  left=10pt, right=10pt, top=6pt, bottom=6pt,
  title=#1,
  fonttitle=\bfseries\color{black},
  attach boxed title to top left={yshift=-2mm, xshift=6mm},
  boxed title style={colback=alertborder, arc=3pt}
}

% caixa de recomendação
\newtcolorbox{recbox}[1][Recomendação]{
  enhanced,
  colback=lightblue!40, colframe=navyblue,
  boxrule=0.8pt, arc=4pt,
  left=10pt, right=10pt, top=6pt, bottom=6pt,
  title=#1,
  fonttitle=\bfseries\color{navyblue},
  attach boxed title to top left={yshift=-2mm, xshift=6mm},
  boxed title style={colback=lightblue, colframe=navyblue, arc=3pt}
}

% ─── Títulos de seções ──────────────────────────────────────────────────────
\usepackage{titlesec}
\titleformat{\section}
  {\color{navyblue}\Large\bfseries}
  {\thesection.}{0.6em}{}
  [\color{navyblue}\titlerule]
\titlespacing{\section}{0pt}{18pt}{8pt}

\titleformat{\subsection}
  {\color{medblue}\large\bfseries}
  {\thesubsection}{0.6em}{}
\titlespacing{\subsection}{0pt}{14pt}{6pt}

\titleformat{\subsubsection}
  {\color{navyblue!80}\normalsize\bfseries}
  {\thesubsubsection}{0.5em}{}
\titlespacing{\subsubsection}{0pt}{10pt}{4pt}

% ─── Hiperlinks ─────────────────────────────────────────────────────────────
\usepackage[
  colorlinks=true,
  linkcolor=navyblue,
  urlcolor=medblue
]{hyperref}

% ─── Listas ─────────────────────────────────────────────────────────────────
\usepackage{enumitem}
\setlist[itemize]{leftmargin=1.5em, itemsep=2pt, topsep=4pt}
\setlist[enumerate]{leftmargin=1.5em, itemsep=2pt, topsep=4pt}

% ─── Outros utilitários ─────────────────────────────────────────────────────
\usepackage{graphicx}
\usepackage{float}
\usepackage{caption}
\captionsetup{font=small, labelfont=bf, textfont=it}

% ============================================================================
\begin{document}
% ============================================================================

% ─── CAPA ───────────────────────────────────────────────────────────────────
\begin{titlepage}
  \pagecolor{navyblue}
  \color{white}
  \centering
  \vspace*{2cm}

  {\large\bfseries SECRETARIA MUNICIPAL DE SAÚDE}\\[4pt]
  {\normalsize MUNICÍPIO DE BORBA --- AMAZONAS}\\[3cm]

  \textcolor{lightblue}{\rule{0.8\textwidth}{1.5pt}}\\[0.6cm]
  {\fontsize{26}{32}\selectfont\bfseries
    Modelagem Matemática e Estatística\\[6pt]
    para a Intendência de Materiais\\[6pt]
    de Limpeza e Higiene}\\[0.6cm]
  \textcolor{lightblue}{\rule{0.8\textwidth}{1.5pt}}\\[1.2cm]

  {\large Rede de Atenção à Saúde do Município de Borba, Amazonas}\\[0.3cm]
  {\normalsize\itshape Exercício 2025--2026}\\[3cm]

  \begin{tabular}{cccc}
    \fbox{\parbox{3cm}{\centering\small\textbf{Documento}\\Estudo Técnico-Científico}} &
    \fbox{\parbox{3cm}{\centering\small\textbf{Orçamento}\\R\$\,2.000.000,00}} &
    \fbox{\parbox{3cm}{\centering\small\textbf{Horizonte}\\12 meses (2025--2026)}} &
    \fbox{\parbox{3cm}{\centering\small\textbf{Versão}\\2.0 --- 2025}}
  \end{tabular}

  \vfill
  {\small Revisado e ampliado com base no estudo original · 2025}
\end{titlepage}
\nopagecolor

% ─── SUMÁRIO ────────────────────────────────────────────────────────────────
\tableofcontents
\clearpage

% ============================================================================
\section{Sumário Executivo}
% ============================================================================

Este estudo fundamenta cientificamente a \textbf{intendência de materiais de limpeza e higiene} da rede de saúde do Município de Borba (AM), aplicando modelagem matemática e estatística adaptada à realidade da Amazônia profunda.

Diferentemente de municípios com acesso rodoviário, Borba opera sob o chamado \textbf{Atrito Logístico Amazônico (ALA)}: combinação de longo prazo de reposição ($TR$), alta variabilidade ($\sigma_{TR}$), impossibilidade de reposição de emergência por via terrestre e sazonalidade do nível do rio Madeira. Os principais resultados são:

\begin{itemize}
  \item Área total de higienização: $A_H \approx 5.150\,\text{m}^2$
  \item Estoque de Segurança ($ES$) mínimo 30\% superior ao de municípios rodoviários
  \item Ponto de Pedido ($PP$) ajustado a $TR = 20$ dias de reposição fluvial
  \item Classificação ABC em 3 categorias com estratégias diferenciadas
  \item Custo por m² majorado pelo frete fluvial Manaus--Borba
\end{itemize}

% ============================================================================
\section{Diagnóstico Estatístico do Ecossistema de Saúde de Borba/AM}
% ============================================================================

O correto dimensionamento pressupõe o mapeamento preciso das variáveis demográficas e de infraestrutura. A Tabela~\ref{tab:variaveis} consolida as variáveis-base adotadas neste estudo.

\vspace{6pt}
\begin{table}[H]
\centering
\caption{Variáveis de diagnóstico do sistema de saúde de Borba/AM}
\label{tab:variaveis}
\small
\rowcolors{2}{rowalt}{white}
\begin{tabular}{>{\bfseries\color{navyblue}}p{1.6cm} p{5.2cm} p{7.2cm}}
  \rowcolor{navyblue}
  \textcolor{white}{\textbf{Variável}} &
  \textcolor{white}{\textbf{Valor}} &
  \textcolor{white}{\textbf{Observação}} \\
  \toprule
  $P$          & 34.869 hab.\ (IBGE 2025)                    & Base para cálculo de demanda per capita e epidemiologia. \\
  $U$          & 10 unidades (9 UBS + 1 Hospital)            & Hospital Vó Mundoca: único serviço de internação do município. \\
  $L$          & 40 leitos operacionais                      & Capacidade do Hospital Vó Mundoca; essencial para o cálculo leito/dia. \\
  $A_C$        & $\approx 2.500\,\text{m}^2$                 & Salas de parto, urgência, UTI e internação --- exigem protocolo de desinfecção concorrente. \\
  $TR$         & 15 a 20 dias                                & Viagem Manaus--Borba em balsa: 11 a 14\,h; frequência semanal limitada e sazonal. \\
  $\sigma_{TR}$& $\pm\,4$ dias                               & Variabilidade estimada pela sazonalidade do nível do Rio Madeira (cheia/seca). \\
  \bottomrule
\end{tabular}
\end{table}

\begin{alertbox}[Variável exclusiva da logística fluvial]
A variável $\sigma_{TR}$ (desvio padrão do tempo de reposição) é exclusiva da logística fluvial e inexiste nos modelos padronizados do Ministério da Saúde. Sua inclusão é \textbf{obrigatória} para o correto dimensionamento do estoque de segurança hospitalar em Borba.
\end{alertbox}

% ============================================================================
\section{Modelagem Matemática da Demanda de Insumos}
% ============================================================================

\subsection{Área Total de Higienização (\texorpdfstring{$A_H$}{AH})}

A área total de higienização $A_H$ representa a superfície total sob responsabilidade da intendência, somando todas as unidades da rede assistencial:

\begin{formulabox}
\[
  A_H \;=\; \sum_{i=1}^{9} \text{UBS}_i \;+\; \text{Area}_{\text{Hosp}}
\]
\end{formulabox}

\noindent Adotando UBS de Porte~I/II ($\approx 350\,\text{m}^2$ cada) e hospital ($\approx 2.000\,\text{m}^2$):

\begin{formulabox}
\[
  A_H \;=\; (9 \times 350) \;+\; 2.000 \;=\; 3.150 \;+\; 2.000 \;=\; \mathbf{5.150\,\text{m}^2}
\]
\end{formulabox}

\begin{recbox}[Recomendação de Levantamento Físico]
  O valor $A_H = 5.150\,\text{m}^2$ é uma estimativa. Recomenda-se o levantamento físico de cada unidade para uso do valor real em licitação.
\end{recbox}

% ─────────────────────────────────────────────────────────────────────────────
\subsection{Subclassificação das Áreas (\texorpdfstring{$A_C$}{AC} e \texorpdfstring{$A_{nc}$}{Anc})}

Para o dimensionamento preciso de saneantes, as áreas são subdivididas em \textbf{críticas} ($A_C$) e \textbf{não críticas} ($A_{nc}$), conforme RDC ANVISA n.º~15/2012 e NR-32:

\vspace{6pt}
\begin{table}[H]
\centering
\caption{Subclassificação das áreas de higienização}
\label{tab:areas}
\small
\rowcolors{2}{rowalt}{white}
\begin{tabular}{p{3.6cm} >{\centering\arraybackslash}p{2.4cm} >{\centering\arraybackslash}p{2.2cm} p{5.4cm}}
  \rowcolor{navyblue}
  \textcolor{white}{\textbf{Classificação}} &
  \textcolor{white}{\textbf{Área (m²)}} &
  \textcolor{white}{\textbf{\% do Total}} &
  \textcolor{white}{\textbf{Exemplos}} \\
  \toprule
  Crítica ($A_C$)           & 2.500 & 48{,}5\% & Salas de cirurgia, UTI, parto, urgência, expurgo. \\
  Não Crítica ($A_{nc}$)    & 2.650 & 51{,}5\% & Corredores, recepção, almoxarifado, vestiários. \\
  \rowcolor{lightblue}
  \textbf{TOTAL ($A_H$)}    & \textbf{5.150} & \textbf{100\%} & \\
  \bottomrule
\end{tabular}
\end{table}

% ─────────────────────────────────────────────────────────────────────────────
\subsection{Consumo Diário de Saneantes (\texorpdfstring{$Q_d$}{Qd})}

O consumo diário $Q_d$ integra as produtividades diferenciadas das áreas críticas e não críticas com o coeficiente de diluição ($\delta$) de cada produto:

\begin{formulabox}
\[
  Q_d \;=\; \left(\frac{A_C}{p_c} \;+\; \frac{A_{nc}}{p_{nc}}\right) \cdot \delta
\]
\end{formulabox}

\noindent onde:
\begin{itemize}
  \item $p_c = 200\,\text{m}^2/\text{dia}$ --- produtividade em área crítica (limpeza terminal)
  \item $p_{nc} = 400\,\text{m}^2/\text{dia}$ --- produtividade em área não crítica (limpeza concorrente)
  \item $\delta$ --- coeficiente de diluição do produto (ex.:~$5\,\text{mL/L}$ para detergente enzimático a 1:200)
\end{itemize}

\begin{alertbox}[Correção de Unidade]
A fórmula original omitia a unidade de resultado de $Q_d$. Com $\delta$ em mL/L e produtividades em m²/dia, o resultado é o \textbf{volume de concentrado por dia (litros)}. Para álcool 70\% (uso puro, $\delta = 1$), o resultado é litros/dia diretamente. Cada produto exige $\delta$ específico conforme ficha técnica do fabricante.
\end{alertbox}

\subsubsection*{Exemplo numérico --- Detergente Enzimático ($\delta = 0{,}005\,\text{L/L}$ a 1:200)}

\begin{formulabox}
\[
  Q_d \;=\; \left(\frac{2500}{200} \;+\; \frac{2650}{400}\right) \cdot 0{,}005
  \;=\; \left(12{,}5 \;+\; 6{,}625\right) \cdot 0{,}005
  \;\approx\; \mathbf{0{,}096\,\text{L}/\text{dia}}
\]
\end{formulabox}

\noindent Consumo mensal estimado: $0{,}096 \times 30 \approx \mathbf{2{,}9\,\text{L de concentrado/mês}}$.

% ============================================================================
\section{Teoremas de Gestão de Estoque e Intendência Amazônica}
% ============================================================================

O \textbf{Atrito Logístico Amazônico (ALA)} combina: longo prazo de reposição ($TR$), alta variabilidade ($\sigma_{TR}$), impossibilidade de reposição de emergência por via rodoviária e sazonalidade do nível do rio. Três equações fundamentais regem a intendência de Borba.

% ─────────────────────────────────────────────────────────────────────────────
\subsection{Estoque de Segurança (\texorpdfstring{$ES$}{ES})}

O estoque de segurança protege o abastecimento durante flutuações imprevisíveis da demanda e do prazo de entrega. A fórmula estocástica completa é:

\begin{formulabox}
\[
  ES \;=\; Z \cdot \sqrt{TR \cdot \sigma_d^2 \;+\; \bar{d}^{\,2} \cdot \sigma_{TR}^2}
\]
\end{formulabox}

\noindent onde:
\begin{itemize}
  \item $Z$ --- fator de nível de serviço (distribuição normal padrão)
  \item $\sigma_d$ --- desvio padrão da demanda diária
  \item $\bar{d}$ --- demanda média diária
  \item $TR$ --- tempo de reposição médio (dias)
  \item $\sigma_{TR}$ --- desvio padrão do tempo de reposição fluvial ($\approx 4$ dias para Borba)
\end{itemize}

\begin{alertbox}[Correção Crítica do Nível de Serviço]
O documento original indicava $Z = 3{,}4$ para 99{,}9\% de disponibilidade. Esta associação é \textbf{incorreta}. O valor correto para 99{,}9\% de nível de serviço (distribuição normal padrão) é $Z = 3{,}09$.

O valor $Z = 3{,}4$ é oriundo do conceito \textbf{Seis Sigma} e refere-se a 3{,}4 defeitos por milhão de oportunidades --- métrica de qualidade de processo, \emph{não} de disponibilidade de estoque.
\end{alertbox}

\begin{table}[H]
\centering
\caption{Valores de $Z$ para dimensionamento do Estoque de Segurança}
\label{tab:z}
\small
\rowcolors{2}{rowalt}{white}
\begin{tabular}{>{\centering\arraybackslash}p{2.8cm} >{\centering\arraybackslash}p{2.2cm} p{8cm}}
  \rowcolor{navyblue}
  \textcolor{white}{\textbf{Nível de Serviço}} &
  \textcolor{white}{\textbf{Valor de $Z$}} &
  \textcolor{white}{\textbf{Recomendação para Borba/AM}} \\
  \toprule
  90{,}0\%            & 1{,}28 & Materiais de baixa criticidade (Classe C). \\
  95{,}0\%            & 1{,}65 & Materiais de criticidade média (Classe B). \\
  99{,}0\%            & 2{,}33 & Piso mínimo para saneantes Classe A em contexto amazônico. \\
  \rowcolor{lightblue}
  \textbf{99{,}9\% \checkmark} & \textbf{3{,}09} & \textbf{RECOMENDADO}: itens críticos do hospital (álcool, hipoclorito, enzimático). \\
  \bottomrule
\end{tabular}
\end{table}

% ─────────────────────────────────────────────────────────────────────────────
\subsection{Ponto de Pedido (\texorpdfstring{$PP$}{PP})}

O Ponto de Pedido define o nível de estoque no qual o pedido deve ser emitido para evitar desabastecimento:

\begin{formulabox}
\[
  PP \;=\; \left(CMM \cdot TR_{\text{meses}}\right) \;+\; ES
\]
\end{formulabox}

\noindent onde:
\begin{itemize}
  \item $CMM$ --- Consumo Médio Mensal
  \item $TR_{\text{meses}}$ --- Tempo de Reposição em meses (Borba: $20\,\text{dias} / 30 \approx 0{,}667\,\text{meses}$)
  \item $ES$ --- Estoque de Segurança calculado pela equação anterior
\end{itemize}

\begin{recbox}[Ajuste Sazonal do $TR$]
Durante a estiagem do Rio Madeira (\textbf{julho--outubro}), o TR deve ser majorado em 20\%, passando de $0{,}667$ para até $0{,}800$ meses, pois o calado mínimo pode atrasar ou impedir balsas de grande porte.
\end{recbox}

% ─────────────────────────────────────────────────────────────────────────────
\subsection{Lote Econômico de Compra --- LEC (\textit{EOQ})}

O LEC (\textit{Economic Order Quantity}) determina a quantidade ótima de pedido que minimiza o custo total de estoque (pedido + armazenagem), especialmente relevante dado o frete fluvial fixo por viagem:

\begin{formulabox}
\[
  \text{LEC} \;=\; \sqrt{\frac{2 \cdot D \cdot C_p}{C_a \cdot P_{\text{unit}}}}
\]
\end{formulabox}

\noindent onde:
\begin{itemize}
  \item $D$ --- demanda anual (unidades)
  \item $C_p$ --- custo por pedido (frete fluvial fixo Manaus--Borba: R\$\,800--1.500 por remessa)
  \item $C_a$ --- taxa de carregamento anual do estoque (20--25\% do valor do item)
  \item $P_{\text{unit}}$ --- preço unitário do item
\end{itemize}

Em Borba, $C_p$ deve incorporar: custo de emissão do empenho + frete fluvial + seguro da carga + perdas estimadas por umidade/temperatura no transporte fluvial (1--3\% para produtos químicos não climatizados).

% ─────────────────────────────────────────────────────────────────────────────
\subsection{Necessidade Real de Aquisição (\texorpdfstring{$N_R$}{NR}) --- Ciclo de 12 Meses}

A quantidade total a ser licitada para o ciclo anual é:

\begin{formulabox}
\[
  N_R \;=\; (CMM \cdot 12) \;+\; ES \;-\; E_{\text{atual}}
\]
\end{formulabox}

Esta equação garante que o processo licitatório cubra o consumo anual projetado, mais a constituição do estoque de segurança, deduzindo o estoque atual para evitar superdimensionamento.

% ============================================================================
\section{Classificação e Curva ABC dos Materiais}
% ============================================================================

A análise ABC prioriza a alocação dos R\$\,2.000.000,00 disponíveis, combinando dois critérios: \textbf{valor financeiro} e \textbf{criticidade biológica/sanitária}.

\begin{table}[H]
\centering
\caption{Curva ABC --- Classificação dos materiais de limpeza e higiene}
\label{tab:abc}
\small
\rowcolors{2}{rowalt}{white}
\begin{tabular}{>{\centering\arraybackslash\bfseries}p{1cm} p{4.8cm} >{\centering\arraybackslash}p{2cm} >{\centering\arraybackslash}p{2cm} p{3.2cm}}
  \rowcolor{navyblue}
  \textcolor{white}{\textbf{Classe}} &
  \textcolor{white}{\textbf{Principais Insumos}} &
  \textcolor{white}{\textbf{\% Valor/Risco}} &
  \textcolor{white}{\textbf{$ES$ Recomendado}} &
  \textcolor{white}{\textbf{Estratégia}} \\
  \toprule
  A & Álcool 70\%, Det.\ Enzimático, Hipoclorito 1\%, Glutaraldeído 2\%, Papel Toalha Hospitalar, EPI de limpeza & 80\% & $Z = 3{,}09$ (99{,}9\%) & Controle diário; licitação própria; emissão automática de $PP$. \\
  B & Sacos p/\ Lixo Infectante, Sabonete Antisséptico, Desinfetante Quaternário, Álcool Isopropílico & 15\% & $Z = 2{,}33$ (99{,}0\%) & Revisão quinzenal; compras semestrais. \\
  C & Vassouras, Rodos, Panos, Baldes, Esponjas, Detergente neutro comum & 5\% & $Z = 1{,}65$ (95{,}0\%) & Estoque de massa; compra anual única. \\
  \bottomrule
\end{tabular}
\end{table}

\begin{alertbox}[Alerta Regulatório]
Todos os saneantes Classe A devem possuir \textbf{Registro ANVISA válido}, conforme Lei n.º~6.360/1976 e RDC n.º~59/2021. A aquisição sem registro configura irregularidade passível de TCU e denúncia ao MP. O edital licitatório deve exigir o número de registro como habilitação técnica.
\end{alertbox}

% ============================================================================
\section{Parâmetros de Consumo Hospitalar --- Hospital Vó Mundoca}
% ============================================================================

\subsection{Consumo de Álcool 70\%}

\begin{itemize}
  \item Higiene das mãos (IRAS): $3\,\text{mL/evento}$; média de 20 eventos/profissional/turno $\Rightarrow 60\,\text{mL/profissional/dia}$
  \item Desinfecção de superfícies críticas: $0{,}5\,\text{L/leito/dia}$ (PNSP/MS)
  \item Desinfecção de equipamentos: $0{,}2\,\text{L/leito/dia}$
\end{itemize}

\begin{formulabox}
\[
  Q_{\text{álcool}} \;=\; 0{,}7\,\text{L/leito/dia} \;\times\; 40\,\text{leitos}
  \;=\; 28\,\text{L/dia} \;\approx\; \mathbf{840\,\text{L/mês}}
\]
\end{formulabox}

\subsection{Consumo de Papel Toalha}

\begin{formulabox}
\[
  Q_{\text{papel}} \;=\; (N_{\text{pacientes}} + N_{\text{profissionais}})
  \;\times\; F_{\text{lavagens}} \;\times\; F_{\text{folhas}}
\]
\end{formulabox}

\noindent Com $F_{\text{lavagens}} = 6\,\text{lavagens/pessoa/dia}$ (OMS) e $F_{\text{folhas}} = 2\,\text{folhas/lavagem}$:

\begin{formulabox}
\[
  Q_{\text{papel}} \;=\; (40 + 60) \;\times\; 6 \;\times\; 2
  \;=\; 1.200\,\text{folhas/dia} \;\approx\; \mathbf{36.000\,\text{folhas/mês}}
\]
\end{formulabox}

\subsection{Consumo de Hipoclorito de Sódio 1\%}

\begin{itemize}
  \item Limpeza terminal (semanal): $0{,}5\,\text{L/m}^2 \times 2.500\,\text{m}^2 = 1.250\,\text{L/semana}$
  \item Desinfecção concorrente diária (0{,}5\%): $0{,}2\,\text{L/m}^2/\text{dia} \times 2.500\,\text{m}^2 = 500\,\text{L/dia}$
\end{itemize}

\begin{alertbox}[Atenção à Degradação Química]
O hipoclorito perde \textbf{20--30\%} da concentração ativa em ambiente de alta temperatura e umidade (condição típica da Amazônia). Recomenda-se embalagens de 5\,L com validade de 30 dias após abertura, em vez de bombonas de 20\,L.
\end{alertbox}

% ============================================================================
\section{Matriz de Custo por Metro Quadrado (\texorpdfstring{$C_{m^2}$}{Cm2})}
% ============================================================================

O custo real de manutenção da salubridade em Borba supera a média nacional em razão do frete fluvial:

\begin{formulabox}
\[
  C_{m^2} \;=\; \frac{\displaystyle\sum \bigl(P_{\text{Manaus}} \;+\; F_{\text{Borba}} \;+\; \text{Perdas}\bigr) \cdot Q}{A_H}
\]
\end{formulabox}

\noindent onde:
\begin{itemize}
  \item $P_{\text{Manaus}}$ --- preço de aquisição no mercado atacadista da capital
  \item $F_{\text{Borba}}$ --- frete fluvial médio por unidade (estimado em 8--15\% do valor CIF Manaus)
  \item $\text{Perdas}$ --- perdas em transporte e armazenagem (1--3\% para líquidos; 0{,}5\% para sólidos)
\end{itemize}

\begin{table}[H]
\centering
\caption{Comparativo de custo real CIF Borba vs.\ FOB Manaus}
\label{tab:custo}
\small
\rowcolors{2}{rowalt}{white}
\begin{tabular}{p{4.2cm} >{\centering\arraybackslash}p{2.6cm} >{\centering\arraybackslash}p{2.6cm} >{\centering\arraybackslash}p{2.6cm} >{\centering\arraybackslash}p{1.8cm}}
  \rowcolor{navyblue}
  \textcolor{white}{\textbf{Produto (Classe A)}} &
  \textcolor{white}{\textbf{Preço Manaus (R\$/L)}} &
  \textcolor{white}{\textbf{Frete Estimado}} &
  \textcolor{white}{\textbf{Custo Real (R\$/L)}} &
  \textcolor{white}{\textbf{Majoração}} \\
  \toprule
  Álcool 70\% (5\,L)         & R\$~12{,}50 & R\$~1{,}63 (13\%) & R\$~14{,}13 & +13\% \\
  Det.\ Enzimático (1\,L)   & R\$~38{,}00 & R\$~4{,}94 (13\%) & R\$~42{,}94 & +13\% \\
  Hipoclorito 1\% (20\,L)   & R\$~4{,}80  & R\$~0{,}72 (15\%) & R\$~5{,}52  & +15\% \\
  \bottomrule
\end{tabular}
\end{table}

\begin{recbox}[Preço de Referência para Licitação]
Recomenda-se que o processo licitatório utilize o \textbf{preço de referência CIF Borba} (não FOB Manaus) para evitar subdimensionamento orçamentário.
\end{recbox}

% ============================================================================
\section{Análise de Sazonalidade e Gestão de Risco}
% ============================================================================

O ciclo hidrológico do Rio Madeira é o principal fator de risco logístico. A intendência deve ser calibrada por \textbf{três períodos operacionais distintos}:

\begin{table}[H]
\centering
\caption{Impacto sazonal do Rio Madeira na logística de suprimentos}
\label{tab:sazonalidade}
\small
\rowcolors{2}{rowalt}{white}
\begin{tabular}{>{\bfseries}p{2.4cm} >{\centering\arraybackslash}p{2cm} p{4.8cm} p{4.4cm}}
  \rowcolor{navyblue}
  \textcolor{white}{\textbf{Período}} &
  \textcolor{white}{\textbf{Meses}} &
  \textcolor{white}{\textbf{Impacto Logístico}} &
  \textcolor{white}{\textbf{Ajuste Operacional}} \\
  \toprule
  Cheia         & Jan -- Jun & Fluxo regular; $TR = 15$--17 dias. Risco de enchente para almoxarifados ribeirinhos. & Estoque padrão. Armazenamento em prateleiras $\geq 50\,\text{cm}$ do solo. \\
  Seca (CRÍTICO)& Jul -- Out & Balsas de maior calado suspendem rotas. $TR$ pode atingir 25--30 dias. Lancha de emergência custa 3$\times$ mais. & $TR = 30$ dias no $PP$. Pedido com 45 dias de antecedência. Elevar $ES$ em 40\%. \\
  Transição     & Nov -- Dez & Rio em recuperação; frequência irregular. $\sigma_{TR}$ elevado. & Usar $\sigma_{TR} = 6$ dias. Antecipar pedidos anuais. \\
  \bottomrule
\end{tabular}
\end{table}

\begin{recbox}[Estoque de Reserva Estratégica (ERE)]
Para itens Classe A, recomenda-se a definição de um \textbf{ERE de 60 dias de consumo}, armazenado fisicamente no Hospital Vó Mundoca, de uso exclusivo em caso de desabastecimento comprovado e autorizado pelo gestor municipal de saúde.
\end{recbox}

% ============================================================================
\section{Indicadores de Desempenho da Intendência (KPIs)}
% ============================================================================

Os seguintes indicadores devem ser monitorados mensalmente pela Comissão de Farmácia e Terapêutica:

\begin{table}[H]
\centering
\caption{KPIs de desempenho da intendência de Borba/AM}
\label{tab:kpis}
\small
\rowcolors{2}{rowalt}{white}
\begin{tabular}{p{4.2cm} p{3.8cm} >{\centering\arraybackslash}p{2.4cm} >{\centering\arraybackslash}p{2.4cm}}
  \rowcolor{navyblue}
  \textcolor{white}{\textbf{Indicador (KPI)}} &
  \textcolor{white}{\textbf{Fórmula}} &
  \textcolor{white}{\textbf{Meta Borba/AM}} &
  \textcolor{white}{\textbf{Periodicidade}} \\
  \toprule
  Taxa de Ruptura de Estoque & $\text{Dias em falta} / \text{Dias do período}$ & $< 0{,}1\%$ & Mensal \\
  Custo por m² Higienizado & $\text{Custo total} / A_H$ & $\leq \text{R\$}\,8{,}50/\text{m}^2$ & Trimestral \\
  Aderência ao $PP$ (Classe A) & $\text{Pedidos no prazo} / \text{Total de pedidos}$ & $> 95\%$ & Mensal \\
  Índice de Perda por Validade & $\text{Valor descartado} / \text{Valor adquirido}$ & $< 2\%$ & Semestral \\
  Cobertura do $ES$ & $E_{\text{atual}} / ES_{\text{calculado}}$ & $> 100\%$ & Quinzenal \\
  \bottomrule
\end{tabular}
\end{table}

% ============================================================================
\section{Conclusão e Recomendações Finais}
% ============================================================================

A intendência de materiais de limpeza e higiene da rede de saúde de Borba/AM não pode ser tratada com os parâmetros-padrão nacionais. O \textbf{Atrito Logístico Amazônico} --- expresso matematicamente pela variável $\sigma_{TR}$ --- exige:

\begin{enumerate}
  \item Estoque de Segurança com $Z = 3{,}09$ ($99{,}9\%$), \emph{não} $Z = 3{,}4$ (erro conceitual do modelo Seis Sigma)
  \item $PP$ ajustado sazonalmente: $TR = 20$ dias na cheia, $TR = 30$ dias na seca do Rio Madeira
  \item LEC que internalize $C_p$ elevado (frete fluvial), favorecendo lotes maiores
  \item Curva ABC com critério combinado financeiro + criticidade sanitária ANVISA
  \item $N_R$ com preço de referência CIF Borba, não FOB Manaus
  \item ERE de 60 dias para itens Classe A durante o período de seca
  \item Monitoramento mensal dos 5 KPIs propostos com relatório ao Conselho Municipal de Saúde
\end{enumerate}

\bigskip
A aplicação rigorosa deste modelo garante que os R\$\,2.000.000,00 disponíveis sejam alocados de forma eficiente e segura, protegendo os \textbf{34.869 habitantes de Borba} contra surtos infecciosos e assegurando a plena operação do \textbf{Hospital Vó Mundoca} durante todo o ciclo hidrológico do Rio Madeira.

\bigskip
\noindent\textcolor{navyblue!50}{\rule{\linewidth}{0.5pt}}

\noindent{\small\textit{Referências técnicas: ANVISA RDC n.º\,15/2012; ANVISA RDC n.º\,59/2021; NR-32 (MTE); PNSP/MS 2013; Tabela $Z$ da Distribuição Normal Padrão (Abramowitz \& Stegun, 1972); Ballou, R.\,H. --- Gerenciamento da Cadeia de Suprimentos, 5.ª ed.\ (2006).}}

\end{document}
